\section{Выводы}
В ходе работы было подтверждено, что обмен между независимыми процессами удобно представлять как передачу данных по отдельным направлениям. Разделение канала передачи чисел и канала ответных сообщений позволило исключить неоднозначность при обмене: родительский процесс отправляет очередное значение только в одном направлении, а результат его обработки получает в другом.

Также было установлено, что корректность взаимодействия процессов зависит не только от передачи данных, но и от управления временем жизни ресурсов. Закрытие неиспользуемых концов каналов необходимо для корректного обнаружения конца потока данных и предотвращения зависания процессов.

Реализованная программа показала, что одну и ту же схему межпроцессного взаимодействия можно сохранить для разных операционных систем, несмотря на различия их API. В Linux и Windows отличаются способы создания процесса и работы с дескрипторами, однако общая модель обмена данными через каналы остаётся одинаковой.
\pagebreak
